\clearpage
\VerdierPhysicalPageAnchor{0064}%
\VerdierSemanticAnchor{ch1-1-6-6-functor-characterization}
命題 (1.5.2) により、$C$ から $D$ への関手を指定することは、関手

\VerdierSemanticAnchor{ch1-1-6-6-eq2}
\begin{equation}
  (X_1,\ldots,X_n)\longmapsto F(X_1,\ldots,X_n)
  \tag{1.6.6.2}
\end{equation}

\VerdierSemanticAnchor{ch1-1-6-6-transition-isomorphisms}
と関手の同型

\VerdierSemanticAnchor{ch1-1-6-6-eq3}
\begin{equation}
  p_i\colon
  F(X_1,\ldots,X_i[1],\ldots,X_n)
  \xrightarrow{\ \sim\ }
  F(X_1,\ldots,X_n)[1]
  \tag{1.6.6.3}
\end{equation}

\VerdierSemanticAnchor{ch1-1-6-6-anticommutativity}
を指定し、図式 (1.5.1.2)

\VerdierSemanticAnchor{ch1-1-6-6-anticommutativity-diagram}
\begin{equation*}
  \resizebox{\textwidth}{!}{%
    $\begin{tikzcd}[ampersand replacement=\&,column sep=2.8em,row sep=3.8em]
      F(X_1,\ldots,X_i[1],\ldots,X_j[1],\ldots,X_n)
        \arrow[r,"{p_i}"]
        \arrow[d,"{p_j}"]
        \& F(X_1,\ldots,X_i,\ldots,X_j[1],\ldots,X_n)[1]
          \arrow[d,"{p_j[1]}"] \\
      F(X_1,\ldots,X_i[1],\ldots,X_j,\ldots,X_n)[1]
        \arrow[r,"{p_i[1]}"]
        \& F(X_1,\ldots,X_i,\ldots,X_j,\ldots,X_n)[2]
    \end{tikzcd}$%
  }
\end{equation*}

が $i<j$ に対して反可換となるようにすることに等しい。（図式 (1.5.1.2) を
簡単にするため、$D$ の遷移同型を恒等射と仮定した。本書では常にそうである。）

\VerdierSemanticAnchor{ch1-1-6-6-morphism-characterization}
同様に命題 (1.5.10) により、ここで考える関手の射
$u\colon F\to F'$ は次の条件に従う。すなわち下の図式は可換である：

\VerdierSemanticAnchor{ch1-1-6-6-eq4}
\begin{equation}
  \resizebox{0.92\textwidth}{!}{%
    $\begin{tikzcd}[ampersand replacement=\&,column sep=5.0em,row sep=3.8em]
      F(X_1,\ldots,X_i[1],\ldots,X_n)
        \arrow[r,"u"]
        \arrow[d,"{p_i}"]
        \& F'(X_1,\ldots,X_i[1],\ldots,X_n)
          \arrow[d,"{p'_i}"] \\
      F(X_1,\ldots,X_i,\ldots,X_n)[1]
        \arrow[r,"{u[1]}"]
        \& F'(X_1,\ldots,X_i,\ldots,X_n)[1]
    \end{tikzcd}$%
  }
  \tag{1.6.6.4}
\end{equation}

\VerdierSemanticAnchor{ch1-1-6-6-underlying-convention}
文脈上混同のおそれがないとき、$\mathbb Z$-圏をその基礎圏で表し、捩られた
$\mathbb Z^n$-関手をその基礎関手で表す。
